\chapter{2013年江苏卷}

\section{填空题}

\begin{problem}
函数 \(y=3\sin\left(2x+\frac{\pi}{4}\right)\) 的最小正周期为\fillinblank{}.
\end{problem}

\begin{answer}
\(\pi\)
\end{answer}

\begin{solution}
正弦函数的最小正周期为
\(T=\dfrac{2\pi}{2}=\pi\).
\end{solution}

\begin{problem}
设 \(z=(2-\i)^2\)（\(\i\) 为虚数单位），则复数 \(z\) 的模为\fillinblank{}.
\end{problem}

\begin{answer}
5
\end{answer}

\begin{solution}
先计算 \(z=(2-\i)^2=4-4\i+\i^2=3-4\i\). 因此
\(|z|=\sqrt{3^2+(-4)^2}=5\).
\end{solution}

\begin{problem}
双曲线 \(\frac{x^2}{16}-\frac{y^2}{9}=1\) 的两条渐近线的方程为\fillinblank{}.
\end{problem}

\begin{answer}
\(y=\pm\frac{3}{4}x\)
\end{answer}

\begin{solution}
双曲线的半实轴长为 \(a=4\)，半虚轴长为 \(b=3\)，其渐近线方程为
\(y=\pm\dfrac{b}{a}x=\pm\dfrac{3}{4}x\).
\end{solution}

\begin{problem}
集合 \(\{-1,0,1\}\) 共有\fillinblank{} 个子集.
\end{problem}

\begin{answer}
8
\end{answer}

\begin{solution}
集合中有 \(3\) 个元素，每个元素都有“取”或“不取”两种选择，所以子集个数为
\(2^3=8\).
\end{solution}

\begin{problem}
如图是一个算法的流程图，则输出的 \(n\) 的值是\fillinblank{}.

\bitmapfigure[max width=0.34\linewidth]{img/2013/jiangsu/q5_fig1.png}
\end{problem}

\begin{answer}
3
\end{answer}

\begin{solution}
初始时 \(n=1\)，\(a=2\). 第一次循环后 \(n=2\)，\(a=3\times2+2=8\)；第二次循环后 \(n=3\)，\(a=3\times8+2=26\). 此时 \(a<20\) 不成立，停止循环并输出 \(n=3\).
\end{solution}

\begin{problem}
抽样统计甲，乙两位射击运动员的 \(5\) 次训练成绩（单位：环），结果如下：

\begin{center}
{\renewcommand{\arraystretch}{1.15}
\begin{tabular}{|c|c|c|c|c|c|}
\hline
运动员 & 第1次 & 第2次 & 第3次 & 第4次 & 第5次 \\
\hline
甲 & 87 & 91 & 90 & 89 & 93 \\
\hline
乙 & 89 & 90 & 91 & 88 & 92 \\
\hline
\end{tabular}}
\end{center}

则成绩较为稳定（方差较小）的那位运动员成绩的方差为\fillinblank{}.
\end{problem}

\begin{answer}
2
\end{answer}

\begin{solution}
甲、乙两人的平均成绩均为 \(90\). 甲的方差为
\(\dfrac{(-3)^2+1^2+0^2+(-1)^2+3^2}{5}=4\)，乙的方差为
\(\dfrac{(-1)^2+0^2+1^2+(-2)^2+2^2}{5}=2\).
因此乙的成绩较为稳定，所求方差为 \(2\).
\end{solution}

\begin{problem}
现有某类病毒记作 \(X_mY_n\)，其中正整数 \(m\)，\(n\)（\(m\leqslant 7\)，\(n\leqslant 9\)）可以任意选取，则 \(m\)，\(n\) 都取到奇数的概率为\fillinblank{}.
\end{problem}

\begin{answer}
\(\frac{20}{63}\)
\end{answer}

\begin{solution}
\(m\) 可取 \(1,2,\ldots,7\)，其中奇数有 \(4\) 个；\(n\) 可取 \(1,2,\ldots,9\)，其中奇数有 \(5\) 个. 所有等可能取法共有 \(7\times9=63\) 种，同时取奇数的取法有 \(4\times5=20\) 种，所以所求概率为 \(\dfrac{20}{63}\).
\end{solution}

\begin{problem}
如图，在三棱柱 \(A_1B_1C_1-ABC\) 中，\(D\)，\(E\)，\(F\) 分别是 \(AB\)，\(AC\)，\(AA_1\) 的中点. 设三棱锥 \(F-ADE\) 的体积为 \(V_1\)，三棱柱 \(A_1B_1C_1-ABC\) 的体积为 \(V_2\)，则 \(V_1:V_2=\fillinblank{}\).

\bitmapfigure[max width=0.42\linewidth]{img/2013/jiangsu/q8_fig1.png}
\end{problem}

\begin{answer}
\(1:24\)
\end{answer}

\begin{solution}
设三棱柱的底面 \(\triangle ABC\) 面积为 \(S\)，高为 \(h\). 因为 \(D\)，\(E\) 分别是 \(AB\)，\(AC\) 的中点，所以 \([\triangle ADE]=\dfrac14S\). 又因为 \(F\) 是 \(AA_1\) 的中点，点 \(F\) 到平面 \(ABC\) 的距离为 \(\dfrac h2\). 因此
\(V_1=\dfrac13\cdot\dfrac14S\cdot\dfrac h2=\dfrac1{24}Sh\)，而 \(V_2=Sh\). 故 \(V_1:V_2=1:24\).
\end{solution}

\begin{problem}
抛物线 \(y=x^2\) 在 \(x=1\) 处的切线与两坐标轴围成的三角形区域为 \(D\)（包含三角形内部与边界）. 若点 \(P(x,y)\) 是区域 \(D\) 内的任意一点，则 \(x+2y\) 的取值范围是\fillinblank{}.
\end{problem}

\begin{answer}
\(\left[-2,\frac{1}{2}\right]\)
\end{answer}

\begin{solution}
函数 \(y=x^2\) 在 \(x=1\) 处的切线斜率为 \(2\)，切线方程为 \(y=2x-1\). 因而三角形区域 \(D\) 的三个顶点为 \((0,0)\)，\(\left(\dfrac12,0\right)\)，\((0,-1)\). 由于 \(x+2y\) 是线性函数，其在三角形区域上的最大值和最小值分别在顶点处取得. 三个顶点处的函数值依次为 \(0\)，\(\dfrac12\)，\(-2\)，所以取值范围为 \(\left[-2,\dfrac12\right]\).
\end{solution}

\begin{problem}
设 \(D\)，\(E\) 分别是 \(\triangle ABC\) 的边 \(AB\)，\(BC\) 上的点，\(AD=\frac{1}{2}AB\)，\(BE=\frac{2}{3}BC\). 若 \(\overrightarrow{DE}=\lambda_1\overrightarrow{AB}+\lambda_2\overrightarrow{AC}\)（\(\lambda_1\)，\(\lambda_2\) 为实数），则 \(\lambda_1+\lambda_2\) 的值为\fillinblank{}.
\end{problem}

\begin{answer}
\(\frac{1}{2}\)
\end{answer}

\begin{solution}
取 \(A\) 为原点. 则 \(\overrightarrow{AD}=\dfrac12\overrightarrow{AB}\)，且
\(\overrightarrow{AE}=\overrightarrow{AB}+\dfrac23\overrightarrow{BC}=\dfrac13\overrightarrow{AB}+\dfrac23\overrightarrow{AC}\). 因此
\(\overrightarrow{DE}=\overrightarrow{AE}-\overrightarrow{AD}=-\dfrac16\overrightarrow{AB}+\dfrac23\overrightarrow{AC}\)，从而 \(\lambda_1=-\dfrac16\)，\(\lambda_2=\dfrac23\)，故 \(\lambda_1+\lambda_2=\dfrac12\).
\end{solution}

\begin{problem}
已知 \(f(x)\) 是定义在 \(\R\) 上的奇函数. 当 \(x>0\) 时，\(f(x)=x^2-4x\)，则不等式 \(f(x)>x\) 的解集用区间表示为\fillinblank{}.
\end{problem}

\begin{answer}
\((-5,0)\cup(5,+\infty)\)
\end{answer}

\begin{solution}
当 \(x>0\) 时，\(f(x)>x\) 等价于 \(x^2-4x>x\)，即 \(x(x-5)>0\)，所以 \(x>5\). 当 \(x<0\) 时，令 \(t=-x>0\). 由奇函数性质，\(f(x)=-f(-x)=-f(t)=-t^2+4t\)，于是 \(f(x)>x\) 等价于 \(-t^2+4t>-t\)，即 \(t(5-t)>0\)，从而 \(-5<x<0\). 当 \(x=0\) 时不等式不成立. 故解集为 \((-5,0)\cup(5,+\infty)\).
\end{solution}

\begin{problem}
在平面直角坐标系 \(xOy\) 中，椭圆 \(C\) 的标准方程为 \(\frac{x^2}{a^2}+\frac{y^2}{b^2}=1\)（\(a>b>0\)），右焦点为 \(F\)，右准线为 \(l\)，短轴的一个端点为 \(B\)，设原点到直线 \(BF\) 的距离为 \(d_1\)，\(F\) 到 \(l\) 的距离为 \(d_2\). 若 \(d_2=\sqrt6\,d_1\)，则椭圆 \(C\) 的离心率为\fillinblank{}.
\end{problem}

\begin{answer}
\(\frac{\sqrt{3}}{3}\)
\end{answer}

\begin{solution}
令椭圆焦距参数为 \(c=\sqrt{a^2-b^2}\)，则离心率为 \(e=\dfrac ca\). 取短轴端点 \(B=(0,b)\)，右焦点 \(F=(c,0)\). 因为 \(c^2+b^2=a^2\)，所以 \(d_1=\dfrac{bc}{\sqrt{b^2+c^2}}=\dfrac{bc}{a}\). 右准线方程为 \(x=\dfrac{a^2}{c}\)，故 \(d_2=\dfrac{a^2}{c}-c=\dfrac{b^2}{c}\). 由题设 \(\dfrac{d_2}{d_1}=\dfrac{ab}{c^2}=\sqrt6\). 代入 \(b=a\sqrt{1-e^2}\)，得到 \(\dfrac{\sqrt{1-e^2}}{e^2}=\sqrt6\). 两边平方并令 \(u=e^2\)，有 \(6u^2+u-1=0\)，所以 \(u=\dfrac13\)（舍去负根），从而 \(e=\dfrac{\sqrt3}{3}\).
\end{solution}

\begin{problem}
在平面直角坐标系 \(xOy\) 中，设定点 \(A(a,a)\)，\(P\) 是函数 \(y=\frac{1}{x}\)（\(x>0\)）图象上一动点. 若点 \(P\)，\(A\) 之间的最短距离为 \(2\sqrt2\)，则满足条件的实数 \(a\) 的所有值为\fillinblank{}.
\end{problem}

\begin{answer}
\(-1,\sqrt{10}\)
\end{answer}

\begin{solution}
设动点 \(P=\left(x,\dfrac1x\right)\)，其中 \(x>0\). 令 \(u=x+\dfrac1x\)，则 \(u\geqslant2\)；反之，每个 \(u\geqslant2\) 都可由某个 \(x>0\) 取得，且
\[
PA^2=\left(x-a\right)^2+\left(\frac1x-a\right)^2=u^2-2-2au+2a^2=(u-a)^2+a^2-2.
\]
若 \(a\geqslant2\)，上式在 \(u=a\) 时取得最小值，故最短距离的平方为 \(a^2-2\). 由题设 \(a^2-2=8\)，得 \(a=\sqrt{10}\)，该值满足 \(a\geqslant2\). 若 \(a<2\)，上式在 \(u=2\) 时取得最小值，故最短距离的平方为 \(2(a-1)^2\). 由 \(2(a-1)^2=8\) 得 \(a=-1\) 或 \(a=3\)，其中只有 \(a=-1\) 满足 \(a<2\). 因此满足条件的实数为 \(-1,\sqrt{10}\).
\end{solution}

\begin{problem}
在正项等比数列 \(\{a_n\}\) 中，\(a_5=\frac{1}{2}\)，\(a_6+a_7=3\)，则满足 \(a_1+a_2+\cdots+a_n>a_1a_2\cdots a_n\) 的最大正整数 \(n\) 的值为\fillinblank{}.
\end{problem}

\begin{answer}
12
\end{answer}

\begin{solution}
设等比数列公比为 \(q>0\). 由 \(a_5=\dfrac12\) 及 \(a_6+a_7=3\)，得 \(\dfrac12q+\dfrac12q^2=3\)，即 \(q^2+q-6=0\). 因为 \(q>0\)，所以 \(q=2\)，从而 \(a_1=\dfrac{a_5}{q^4}=\dfrac1{32}\).

于是 \(S_n=\dfrac{2^n-1}{32}\)，且 \(a_1a_2\cdots a_n=\left(\dfrac1{32}\right)^n2^{n(n-1)/2}=2^{(n^2-11n)/2}\). 不等式等价于 \(2^n-1>2^{(n^2-11n+10)/2}\). 当 \(2\leqslant n\leqslant12\) 时，\(n^2-13n+10<0\)，所以右端指数小于 \(n\)，进而右端不超过 \(2^{n-1}<2^n-1\)，不等式成立. 当 \(n\geqslant13\) 时，\(n^2-13n+10>0\)，右端指数大于 \(n\)，故右端大于 \(2^n-1\)，不等式不成立. 当 \(n=1\) 时两边相等. 因此最大正整数为 \(12\).
\end{solution}

\section{解答题}

\begin{problem}
已知 \(\bs{a}=(\cos\alpha,\sin\alpha)\)，\(\bs{b}=(\cos\beta,\sin\beta)\)，\(0<\beta<\alpha<\pi\).

\begin{enumerate}
\item 若 \(\left|\bs{a}-\bs{b}\right|=\sqrt2\)，求证：\(\bs{a}\perp\bs{b}\)；
\item 设 \(\bs{c}=(0,1)\)，若 \(\bs{a}+\bs{b}=\bs{c}\)，求 \(\alpha\)，\(\beta\) 的值.
\end{enumerate}
\end{problem}

\begin{solution}
\begin{enumerate}
\item 因为 \(\lvert\bs{a}\rvert=\lvert\bs{b}\rvert=1\)，所以
\(\lvert\bs{a}-\bs{b}\rvert^2=\lvert\bs{a}\rvert^2+\lvert\bs{b}\rvert^2-2\bs{a}\cdot\bs{b}=2-2\bs{a}\cdot\bs{b}\). 由题设 \(\lvert\bs{a}-\bs{b}\rvert=\sqrt2\)，得 \(\bs{a}\cdot\bs{b}=0\)，因此 \(\bs{a}\perp\bs{b}\).
\item 由 \(\bs{a}+\bs{b}=\bs{c}\) 的两个坐标分量，得 \(\cos\alpha+\cos\beta=0\)，\(\sin\alpha+\sin\beta=1\). 由和差化积公式，第一式为
\(2\cos\dfrac{\alpha+\beta}{2}\cos\dfrac{\alpha-\beta}{2}=0\). 由于 \(0<\alpha-\beta<\pi\)，有 \(\cos\dfrac{\alpha-\beta}{2}>0\)，且 \(0<\dfrac{\alpha+\beta}{2}<\pi\)，所以 \(\alpha+\beta=\pi\). 再由第二式得 \(2\sin\dfrac{\alpha+\beta}{2}\cos\dfrac{\alpha-\beta}{2}=1\)，即 \(\cos\dfrac{\alpha-\beta}{2}=\dfrac12\). 由于 \(0<\dfrac{\alpha-\beta}{2}<\dfrac\pi2\)，故 \(\dfrac{\alpha-\beta}{2}=\dfrac\pi3\). 解得 \(\alpha=\dfrac{5\pi}{6}\)，\(\beta=\dfrac\pi6\).
\end{enumerate}
\end{solution}

\begin{problem}
如图，在三棱锥 \(S-ABC\) 中，平面 \(SAB\perp\) 平面 \(SBC\)，\(AB\perp BC\)，\(AS=AB\). 过 \(A\) 作 \(AF\perp SB\)，垂足为 \(F\)，点 \(E\)，\(G\) 分别是棱 \(SA\)，\(SC\) 的中点. 求证：

\begin{enumerate}
\item 平面 \(EFG\parallel\) 平面 \(ABC\)；
\item \(BC\perp SA\).
\end{enumerate}

\bitmapfigure[max width=0.42\linewidth]{img/2013/jiangsu/q16_fig1.png}
\end{problem}

\begin{solution}
\begin{enumerate}
\item 在等腰三角形 \(ASB\) 中，\(AS=AB\)，且 \(AF\perp SB\)，所以高 \(AF\) 同时是中线，\(F\) 为 \(SB\) 的中点. 又因为 \(E\)，\(F\) 分别是 \(SA\)，\(SB\) 的中点，故在三角形 \(SAB\) 中 \(EF\parallel AB\). 在三角形 \(SAC\) 中，\(E\)，\(G\) 分别是 \(SA\)，\(SC\) 的中点，故 \(EG\parallel AC\). 直线 \(EF\)，\(EG\) 分别平行于平面 \(ABC\) 内相交的直线 \(AB\)，\(AC\)，所以平面 \(EFG\parallel\) 平面 \(ABC\).
\item 两平面 \(SAB\)，\(SBC\) 的交线为 \(SB\). 因为 \(AF\) 在平面 \(SAB\) 内且 \(AF\perp SB\)，由二面角的性质知 \(AF\perp\) 平面 \(SBC\)，从而 \(AF\perp BC\). 又已知 \(AB\perp BC\)，而 \(\overrightarrow{AB}=\overrightarrow{AF}+\overrightarrow{FB}\)，所以 \(FB\perp BC\). 因为 \(F\)，\(B\)，\(S\) 共线，故 \(SB\perp BC\). 直线 \(SA\) 位于由 \(AF\) 和 \(SB\) 所确定的平面内，且 \(BC\) 同时垂直于这两条相交直线，因此 \(BC\perp SA\).
\end{enumerate}
\end{solution}

\begin{problem}
如图，在平面直角坐标系 \(xOy\) 中，点 \(A(0,3)\)，直线 \(l:y=2x-4\). 设圆 \(C\) 的半径为 \(1\)，圆心在 \(l\) 上.

\begin{enumerate}
\item 若圆心 \(C\) 也在直线 \(y=x-1\) 上，过点 \(A\) 作圆 \(C\) 的切线，求切线的方程；
\item 若圆 \(C\) 上存在点 \(M\)，使 \(MA=2MO\)，求圆心 \(C\) 的横坐标 \(a\) 的取值范围.
\end{enumerate}

\bitmapfigure[max width=0.42\linewidth]{img/2013/jiangsu/q17_fig1.png}
\end{problem}

\begin{solution}
\begin{enumerate}
\item 设圆心为 \(C=(a,2a-4)\). 由圆心还在直线 \(y=x-1\) 上，得 \(2a-4=a-1\)，所以 \(a=3\)，圆心为 \(C=(3,2)\). 过 \(A(0,3)\) 的非竖直直线可设为 \(y-3=kx\)，即 \(kx-y+3=0\). 圆心到该直线的距离等于半径 \(1\)，故
\(\dfrac{|3k-2+3|}{\sqrt{k^2+1}}=1\). 于是 \((3k+1)^2=k^2+1\)，即 \(8k^2+6k=0\)，所以 \(k=0\) 或 \(k=-\dfrac34\). 直线 \(x=0\) 到圆心的距离为 \(3\)，不是切线，故全部切线为
\(y=3\) 和 \(y-3=-\dfrac34x\)，即 \(y=3\) 或 \(3x+4y-12=0\).
\item 设 \(M=(x,y)\). 条件 \(MA=2MO\) 等价于
\(x^2+(y-3)^2=4(x^2+y^2)\)，整理得 \(x^2+(y+1)^2=4\). 因此 \(M\) 在圆 \(\Gamma\colon x^2+(y+1)^2=4\) 上，该圆的圆心为 \(N=(0,-1)\)，半径为 \(2\). 题设要求半径为 \(1\) 的圆 \(C\) 与 \(\Gamma\) 有公共点，所以
\(1\leqslant CN\leqslant3\). 又 \(C=(a,2a-4)\)，从而
\(CN^2=a^2+(2a-3)^2=5a^2-12a+9\). 下界条件等价于 \(5a^2-12a+8\geqslant0\)，其判别式为 \((-12)^2-4\times5\times8=-16<0\)，恒成立；上界条件为 \(5a^2-12a\leqslant0\)，即 \(a(5a-12)\leqslant0\). 因此 \(a\) 的取值范围为 \(\left[0,\dfrac{12}{5}\right]\).
\end{enumerate}
\end{solution}

\begin{problem}
如图，游客从某旅游景区的景点 \(A\) 处下山至 \(C\) 处有两种路径. 一种是从 \(A\) 沿直线步行到 \(C\)，另一种是先从 \(A\) 沿索道乘缆车到 \(B\)，然后从 \(B\) 沿直线步行到 \(C\). 现有甲，乙两位游客从 \(A\) 处下山，甲沿 \(AC\) 匀速步行，速度为 \(50\mathrm{~m/min}\). 在甲出发 \(2\mathrm{~min}\) 后，乙从 \(A\) 乘缆车到 \(B\)，在 \(B\) 处停留 \(1\mathrm{~min}\) 后，再从 \(B\) 匀速步行到 \(C\). 假设缆车匀速直线运动的速度为 \(130\mathrm{~m/min}\)，山路 \(AC\) 长为 \(1260\mathrm{~m}\)，经测量，\(\cos A=\frac{12}{13}\)，\(\cos C=\frac{3}{5}\).

\begin{enumerate}
\item 求索道 \(AB\) 的长；
\item 问乙出发多少分钟后，乙在缆车上与甲的距离最短？
\item 为使两位游客在 \(C\) 处互相等待的时间不超过 \(3\mathrm{~min}\)，乙步行的速度应控制在什么范围内？
\end{enumerate}

\bitmapfigure[max width=0.42\linewidth]{img/2013/jiangsu/q18_fig1.png}
\end{problem}

\begin{solution}
由 \(\cos A=\dfrac{12}{13}\)，\(\cos C=\dfrac35\)，得 \(\sin A=\dfrac5{13}\)，\(\sin C=\dfrac45\). 又
\(\sin B=\sin(A+C)=\sin A\cos C+\cos A\sin C=\dfrac{63}{65}\). 由正弦定理，\(AC=1260\) 对应角 \(B\)，所以
\(AB=1260\cdot\dfrac{\sin C}{\sin B}=1260\cdot\dfrac{4/5}{63/65}=1040\)，且
\(BC=1260\cdot\dfrac{\sin A}{\sin B}=1260\cdot\dfrac{5/13}{63/65}=500\).

\begin{enumerate}
\item 索道 \(AB\) 的长为 \(1040\mathrm{~m}\).
\item 以 \(A\) 为原点、\(AC\) 所在直线为 \(x\) 轴正方向建立直角坐标系. 由 \(AB=1040\) 及 \(\cos A=\dfrac{12}{13}\)，可取 \(B=(960,400)\). 设 \(t\) 为乙出发后的时间，乙在缆车上时 \(0\leqslant t\leqslant\dfrac{1040}{130}=8\)，故乙的位置为 \((120t,50t)\). 此时甲已行走 \(50(t+2)\mathrm{~m}\)，位置为 \((50t+100,0)\). 两人的距离平方为
\(d^2(t)=(120t-50t-100)^2+(50t)^2=7400t^2-14000t+10000\).
由于这是开口向上的二次函数，其最小值在 \(t=\dfrac{14000}{2\times7400}=\dfrac{35}{37}\) 时取得，且该值属于 \([0,8]\). 所以乙出发 \(\dfrac{35}{37}\mathrm{~min}\) 后两人距离最短.
\item 甲到达 \(C\) 的时刻为 \(\dfrac{1260}{50}=\dfrac{126}{5}\mathrm{~min}\). 乙乘缆车用时 \(8\mathrm{~min}\)，在 \(B\) 停留 \(1\mathrm{~min}\)，所以乙开始步行的时刻为 \(2+8+1=11\mathrm{~min}\). 若乙步行速度为 \(v\mathrm{~m/min}\)，则乙到达 \(C\) 的时刻为 \(11+\dfrac{500}{v}\). 两人在 \(C\) 处互相等待的时间不超过 \(3\mathrm{~min}\)，等价于
\(\left|11+\dfrac{500}{v}-\dfrac{126}{5}\right|\leqslant3\).
因 \(v>0\)，这又等价于 \(\dfrac{56}{5}\leqslant\dfrac{500}{v}\leqslant\dfrac{86}{5}\)，从而
\(\dfrac{1250}{43}\leqslant v\leqslant\dfrac{625}{14}\). 因此乙步行速度应控制在 \(\left[\dfrac{1250}{43},\dfrac{625}{14}\right]\mathrm{~m/min}\) 内.
\end{enumerate}
\end{solution}

\begin{problem}
设 \(\{a_n\}\) 是首项为 \(a\)，公差为 \(d\) 的等差数列（\(d\ne 0\)），\(S_n\) 是其前 \(n\) 项和. 记 \(b_n=\frac{nS_n}{n^2+c}\)，\(n\in\N^*\)，其中 \(c\) 为实数.

\begin{enumerate}
\item 若 \(c=0\)，且 \(b_1\)，\(b_2\)，\(b_4\) 成等比数列，证明：\(S_{nk}=n^2S_k\)（\(k\)，\(n\in\N^*\)）；
\item 若 \(\{b_n\}\) 是等差数列，证明：\(c=0\).
\end{enumerate}
\end{problem}

\begin{solution}
\begin{enumerate}
\item 当 \(c=0\) 时，\(b_n=\dfrac{S_n}{n}=a+\dfrac{n-1}{2}d\). 由 \(b_1\)，\(b_2\)，\(b_4\) 成等比数列，得 \(\left(a+\dfrac d2\right)^2=a\left(a+\dfrac{3d}{2}\right)\). 整理得 \(d(d-2a)=0\). 因为 \(d\ne0\)，所以 \(d=2a\). 于是 \(S_n=\dfrac n2\left[2a+(n-1)d\right]=\dfrac{n^2d}{2}\)，从而 \(S_{nk}=\dfrac{(nk)^2d}{2}=n^2\dfrac{k^2d}{2}=n^2S_k\).
\item 记 \(p=\dfrac d2\ne0\)，\(q=a-\dfrac d2\). 则 \(S_n=pn^2+qn\)，所以 \(b_n=\dfrac{pn^3+qn^2}{n^2+c}\). 若 \(\{b_n\}\) 是等差数列，设其通项为 \(un+v\). 当 \(n\) 趋于无穷时，\(\dfrac{b_n}{n}\to p\)，故 \(u=p\)；又 \(b_n-pn\to q\)，故 \(v=q\). 因而对每个正整数 \(n\)，有 \(b_n=pn+q\). 代入表达式得
\(pn^3+qn^2=(pn+q)(n^2+c)=pn^3+qn^2+c(pn+q)\). 因为 \(p\ne0\)，上式对所有正整数 \(n\) 成立只能有 \(c=0\).
\end{enumerate}
\end{solution}

\begin{problem}
设函数 \(f(x)=\ln x-ax\)，\(g(x)=\e^x-ax\)，其中 \(a\) 为实数.

\begin{enumerate}
\item 若 \(f(x)\) 在 \((1,+\infty)\) 上是单调减函数，且 \(g(x)\) 在 \((1,+\infty)\) 上有最小值，求 \(a\) 的取值范围；
\item 若 \(g(x)\) 在 \((-1,+\infty)\) 上是单调增函数，试求 \(f(x)\) 的零点个数，并证明你的结论.
\end{enumerate}
\end{problem}

\begin{solution}
\begin{enumerate}
\item \(f'(x)=\dfrac1x-a\). 要使 \(f(x)\) 在 \((1,+\infty)\) 上单调减，需有 \(\dfrac1x-a\leqslant0\) 对一切 \(x>1\) 成立，这给出 \(a\geqslant1\). 又 \(g'(x)=\e^x-a\)，且 \(g''(x)=\e^x>0\). 当且仅当 \(a>\e\) 时，方程 \(g'(x)=0\) 的根 \(x=\ln a\) 落在 \((1,+\infty)\) 内，此时 \(g\) 先减后增，在 \(x=\ln a\) 处取得最小值. 因此 \(a\) 的取值范围为 \((\e,+\infty)\).
\item 由 \(g'(x)=\e^x-a\)，\(g(x)\) 在 \((-1,+\infty)\) 上单调增的充要条件为 \(a\leqslant\e^{-1}\). 在此条件下考察 \(f(x)=\ln x-ax\) 的零点.

当 \(a\leqslant0\) 时，\(f'(x)=\dfrac1x-a>0\)，且 \(f(x)\) 从 \(x\to0^+\) 时的 \(-\infty\) 增至 \(x\to+\infty\) 时的 \(+\infty\)，所以有且只有一个零点.

当 \(0<a\leqslant\e^{-1}\) 时，\(f'(x)=0\) 的唯一根为 \(x=\dfrac1a\)，函数在 \(\left(0,\dfrac1a\right)\) 上递增，在 \(\left(\dfrac1a,+\infty\right)\) 上递减. 且
\(f\left(\dfrac1a\right)=\ln\dfrac1a-1=-\ln a-1\). 当 \(0<a<\e^{-1}\) 时，该最大值大于 \(0\)，结合两端极限均为 \(-\infty\)，故有两个零点；当 \(a=\e^{-1}\) 时最大值等于 \(0\)，故有一个零点 \(x=\e\). 所以零点个数为
\(\begin{cases}
1,&a\leqslant0,\\
2,&0<a<\e^{-1},\\
1,&a=\e^{-1}.
\end{cases}\)
\end{enumerate}
\end{solution}

\section{附加题：解答题}

\begin{problem}[A]
如图，\(AB\) 和 \(BC\) 分别与圆 \(O\) 相切于点 \(D\)，\(C\)，\(AC\) 经过圆心 \(O\)，且 \(BC=2OC\). 求证：\(AC=2AD\).

\bitmapfigure[max width=0.42\linewidth]{img/2013/jiangsu/q21_fig1.png}
\end{problem}

\begin{solution}
设 \(OC=r\). 建立直角坐标系，使 \(O=(0,0)\)，\(C=(r,0)\)，由于 \(BC\) 是在 \(C\) 处的切线且 \(BC=2r\)，可取 \(B=(r,2r)\). 设另一切点为 \(D=(x,y)\). 由 \(D\) 在圆上及 \(OD\perp BD\)，有 \(x^2+y^2=r^2\)，\((B-D)\cdot D=0\)，即 \(x+2y=r\). 联立得除 \(C=(r,0)\) 外的另一解为 \(D=\left(-\dfrac{3r}{5},\dfrac{4r}{5}\right)\).

圆在 \(D\) 处的切线方程为 \(xx_D+yy_D=r^2\)，即 \(-3x+4y=5r\). 因为 \(A\)，\(O\)，\(C\) 共线，\(A\) 在 \(x\) 轴上，所以 \(A=\left(-\dfrac{5r}{3},0\right)\). 因此 \(AC=r+\dfrac{5r}{3}=\dfrac{8r}{3}\)，而
\(AD=\sqrt{\left(\dfrac{16r}{15}\right)^2+\left(\dfrac{4r}{5}\right)^2}=\dfrac{4r}{3}\). 故 \(AC=2AD\).
\end{solution}

\reuseproblemnumber
\begin{problem}[B]
已知矩阵 \(A=\begin{bmatrix}-1&0\\0&2\end{bmatrix}\)，\(B=\begin{bmatrix}1&2\\0&6\end{bmatrix}\)，求矩阵 \(A^{-1}B\).
\end{problem}

\begin{solution}
因为 \(A=\begin{bmatrix}-1&0\\0&2\end{bmatrix}\)，所以 \(A^{-1}=\begin{bmatrix}-1&0\\0&\dfrac12\end{bmatrix}\). 因而
\(A^{-1}B=\begin{bmatrix}-1&0\\0&\dfrac12\end{bmatrix}\begin{bmatrix}1&2\\0&6\end{bmatrix}=\begin{bmatrix}-1&-2\\0&3\end{bmatrix}\).
\end{solution}

\reuseproblemnumber
\begin{problem}[C]
在平面直角坐标系 \(xOy\) 中，直线 \(l\) 的参数方程为 \(x=t+1\)，\(y=2t\)（\(t\) 为参数），曲线 \(C\) 的参数方程为 \(x=2\tan^2\theta\)，\(y=2\tan\theta\)（\(\theta\) 为参数）. 试求直线 \(l\) 和曲线 \(C\) 的普通方程，并求出它们的公共点的坐标.
\end{problem}

\begin{solution}
由直线的参数方程消去参数 \(t\)，得 \(t=\dfrac y2\)，所以直线 \(l\) 的普通方程为 \(y=2x-2\). 由曲线的参数方程消去 \(\theta\)，得 \(\tan\theta=\dfrac y2\)，从而 \(x=2\tan^2\theta=\dfrac{y^2}{2}\)，即曲线 \(C\) 的普通方程为 \(y^2=2x\).

联立 \(y=2x-2\) 与 \(x=\dfrac{y^2}{2}\)，得 \(y=y^2-2\)，即 \(y^2-y-2=0\). 因此 \(y=2\) 或 \(y=-1\)，相应的 \(x=\dfrac{y^2}{2}\) 分别为 \(2\) 和 \(\dfrac12\). 所以公共点坐标为 \((2,2)\) 和 \(\left(\dfrac12,-1\right)\).
\end{solution}

\reuseproblemnumber
\begin{problem}[D]
已知 \(a\geqslant b>0\)，求证：\(2a^3-b^3\geqslant 2ab^2-a^2b\).
\end{problem}

\begin{solution}
将不等式左边减去右边并因式分解，得
\(2a^3-b^3-2ab^2+a^2b=(a-b)(2a^2+3ab+b^2)=(a-b)(a+b)(2a+b)\).
因为 \(a\geqslant b>0\)，所以 \(a-b\geqslant0\)，\(a+b>0\)，\(2a+b>0\)，故上述乘积不小于 \(0\)，从而 \(2a^3-b^3\geqslant2ab^2-a^2b\).
\end{solution}

\begin{problem}
如图，在直三棱柱 \(A_1B_1C_1-ABC\) 中，\(AB\perp AC\)，\(AB=AC=2\)，\(A_1A=4\)，点 \(D\) 是 \(BC\) 的中点.

\begin{enumerate}
\item 求异面直线 \(A_1B\) 与 \(C_1D\) 所成角的余弦值；
\item 求平面 \(ADC_1\) 与平面 \(ABA_1\) 所成二面角的正弦值.
\end{enumerate}

\bitmapfigure[max width=0.42\linewidth]{img/2013/jiangsu/q22_fig1.png}
\end{problem}

\begin{solution}
建立空间直角坐标系，取 \(A=(0,0,0)\)，\(B=(2,0,0)\)，\(C=(0,2,0)\)，\(A_1=(0,0,4)\). 则 \(C_1=(0,2,4)\)，且 \(D=(1,1,0)\).
\begin{enumerate}
\item 可取异面直线 \(A_1B\)，\(C_1D\) 的方向向量分别为 \(\overrightarrow{A_1B}=(2,0,-4)\)，\(\overrightarrow{C_1D}=(1,-1,-4)\). 所以
\(\cos\theta=\dfrac{\left|\overrightarrow{A_1B}\cdot\overrightarrow{C_1D}\right|}{\lvert\overrightarrow{A_1B}\rvert\,\lvert\overrightarrow{C_1D}\rvert}=\dfrac{|2+16|}{\sqrt{20}\sqrt{18}}=\dfrac{3\sqrt{10}}{10}\).
\item 平面 \(ADC_1\) 内的两个相交向量可取 \(\overrightarrow{AD}=(1,1,0)\)，\(\overrightarrow{AC_1}=(0,2,4)\)，故其法向量可取
\(\bs{n}_1=\overrightarrow{AD}\times\overrightarrow{AC_1}=(4,-4,2)\). 平面 \(ABA_1\) 的法向量可取 \(\bs{n}_2=(0,1,0)\). 设两平面的锐二面角为 \(\varphi\)，则
\(\sin\varphi=\dfrac{|\bs{n}_1\times\bs{n}_2|}{|\bs{n}_1|\,|\bs{n}_2|}=\dfrac{\sqrt{20}}{6}=\dfrac{\sqrt5}{3}\).
\end{enumerate}
\end{solution}

\begin{problem}
设数列 \(\{a_n\}\)：\(1\)，\(-2\)，\(-2\)，\(3\)，\(3\)，\(3\)，\(-4\)，\(-4\)，\(-4\)，\(-4\)，\(\cdots\)，\(\underbrace{(-1)^{k-1}k,\cdots,(-1)^{k-1}k}_{k\text{ 个}}\)，\(\cdots\). 即当 \(\frac{(k-1)k}{2}<n\leqslant \frac{k(k+1)}{2}\)（\(k\in\N^*\)）时，\(a_n=(-1)^{k-1}k\). 记 \(S_n=a_1+a_2+\cdots+a_n\)（\(n\in\N^*\)）. 对于 \(\ell\in\N^*\)，定义集合 \(P_{\ell}=\{n\mid 1\leqslant n\leqslant \ell,\ S_n\text{ 是 }a_n\text{ 的整数倍}\}\).

\begin{enumerate}
\item 求集合 \(P_{11}\) 中元素的个数；
\item 求集合 \(P_{2000}\) 中元素的个数.
\end{enumerate}
\end{problem}

\begin{solution}
记第 \(k\) 组的末项下标为 \(T_k=\dfrac{k(k+1)}2\). 第 \(k\) 组共有 \(k\) 项，每项均为 \((-1)^{k-1}k\)，所以前 \(k-1\) 组的和为
\(S_{T_{k-1}}=\sum_{j=1}^{k-1}(-1)^{j-1}j^2\).

若 \(k=2p+1\) 为奇数，则
\(S_{T_{k-1}}=1^2-2^2+\cdots-(2p)^2=-p(2p+1)=-\dfrac{k(k-1)}2\). 在第 \(k\) 组内令 \(n=T_{k-1}+r\)，其中 \(1\leqslant r\leqslant k\)，则
\(S_n=-\dfrac{k(k-1)}2+kr=k\left(r-\dfrac{k-1}{2}\right)\)，必为 \(a_n=k\) 的整数倍.

若 \(k=2p\) 为偶数，则
\(S_{T_{k-1}}=1^2-2^2+\cdots+(2p-1)^2=p(2p-1)=\dfrac{k(k-1)}2\). 同样令 \(n=T_{k-1}+r\)，则
\(S_n=\dfrac{k(k-1)}2-kr=-k\left(r-\dfrac{k-1}{2}\right)\). 因为 \(k-1\) 为奇数，\(r-\dfrac{k-1}{2}\) 为半整数，故 \(S_n\) 不是 \(a_n=-k\) 的整数倍.

因此，恰有奇数编号的组中的下标属于集合 \(P_\ell\). 由于 \(T_4=10<T_5=15\)，前 \(11\) 项中包含第 \(1\)，\(3\) 组的全部项及第 \(5\) 组的第 \(1\) 项，故
\(|P_{11}|=1+3+1=5\).

又因为 \(T_{62}=1953<2000\leqslant T_{63}=2016\)，前 \(2000\) 项中第 \(1,3,\ldots,61\) 组全部符合条件，共有
\(1+3+\cdots+61=31^2=961\) 项；第 \(63\) 组的前 \(2000-1953=47\) 项也全部符合条件. 所以
\(|P_{2000}|=961+47=1008\).
\end{solution}
